% Options for packages loaded elsewhere
% Options for packages loaded elsewhere
\PassOptionsToPackage{unicode}{hyperref}
\PassOptionsToPackage{hyphens}{url}
\PassOptionsToPackage{dvipsnames,svgnames,x11names}{xcolor}
%
\documentclass[
  letterpaper,
  DIV=11,
  numbers=noendperiod]{scrartcl}
\usepackage{xcolor}
\usepackage{amsmath,amssymb}
\setcounter{secnumdepth}{-\maxdimen} % remove section numbering
\usepackage{iftex}
\ifPDFTeX
  \usepackage[T1]{fontenc}
  \usepackage[utf8]{inputenc}
  \usepackage{textcomp} % provide euro and other symbols
\else % if luatex or xetex
  \usepackage{unicode-math} % this also loads fontspec
  \defaultfontfeatures{Scale=MatchLowercase}
  \defaultfontfeatures[\rmfamily]{Ligatures=TeX,Scale=1}
\fi
\usepackage{lmodern}
\ifPDFTeX\else
  % xetex/luatex font selection
\fi
% Use upquote if available, for straight quotes in verbatim environments
\IfFileExists{upquote.sty}{\usepackage{upquote}}{}
\IfFileExists{microtype.sty}{% use microtype if available
  \usepackage[]{microtype}
  \UseMicrotypeSet[protrusion]{basicmath} % disable protrusion for tt fonts
}{}
\makeatletter
\@ifundefined{KOMAClassName}{% if non-KOMA class
  \IfFileExists{parskip.sty}{%
    \usepackage{parskip}
  }{% else
    \setlength{\parindent}{0pt}
    \setlength{\parskip}{6pt plus 2pt minus 1pt}}
}{% if KOMA class
  \KOMAoptions{parskip=half}}
\makeatother
% Make \paragraph and \subparagraph free-standing
\makeatletter
\ifx\paragraph\undefined\else
  \let\oldparagraph\paragraph
  \renewcommand{\paragraph}{
    \@ifstar
      \xxxParagraphStar
      \xxxParagraphNoStar
  }
  \newcommand{\xxxParagraphStar}[1]{\oldparagraph*{#1}\mbox{}}
  \newcommand{\xxxParagraphNoStar}[1]{\oldparagraph{#1}\mbox{}}
\fi
\ifx\subparagraph\undefined\else
  \let\oldsubparagraph\subparagraph
  \renewcommand{\subparagraph}{
    \@ifstar
      \xxxSubParagraphStar
      \xxxSubParagraphNoStar
  }
  \newcommand{\xxxSubParagraphStar}[1]{\oldsubparagraph*{#1}\mbox{}}
  \newcommand{\xxxSubParagraphNoStar}[1]{\oldsubparagraph{#1}\mbox{}}
\fi
\makeatother


\usepackage{longtable,booktabs,array}
\usepackage{calc} % for calculating minipage widths
% Correct order of tables after \paragraph or \subparagraph
\usepackage{etoolbox}
\makeatletter
\patchcmd\longtable{\par}{\if@noskipsec\mbox{}\fi\par}{}{}
\makeatother
% Allow footnotes in longtable head/foot
\IfFileExists{footnotehyper.sty}{\usepackage{footnotehyper}}{\usepackage{footnote}}
\makesavenoteenv{longtable}
\usepackage{graphicx}
\makeatletter
\newsavebox\pandoc@box
\newcommand*\pandocbounded[1]{% scales image to fit in text height/width
  \sbox\pandoc@box{#1}%
  \Gscale@div\@tempa{\textheight}{\dimexpr\ht\pandoc@box+\dp\pandoc@box\relax}%
  \Gscale@div\@tempb{\linewidth}{\wd\pandoc@box}%
  \ifdim\@tempb\p@<\@tempa\p@\let\@tempa\@tempb\fi% select the smaller of both
  \ifdim\@tempa\p@<\p@\scalebox{\@tempa}{\usebox\pandoc@box}%
  \else\usebox{\pandoc@box}%
  \fi%
}
% Set default figure placement to htbp
\def\fps@figure{htbp}
\makeatother





\setlength{\emergencystretch}{3em} % prevent overfull lines

\providecommand{\tightlist}{%
  \setlength{\itemsep}{0pt}\setlength{\parskip}{0pt}}



 


\KOMAoption{captions}{tableheading}
\makeatletter
\@ifpackageloaded{caption}{}{\usepackage{caption}}
\AtBeginDocument{%
\ifdefined\contentsname
  \renewcommand*\contentsname{Table of contents}
\else
  \newcommand\contentsname{Table of contents}
\fi
\ifdefined\listfigurename
  \renewcommand*\listfigurename{List of Figures}
\else
  \newcommand\listfigurename{List of Figures}
\fi
\ifdefined\listtablename
  \renewcommand*\listtablename{List of Tables}
\else
  \newcommand\listtablename{List of Tables}
\fi
\ifdefined\figurename
  \renewcommand*\figurename{Figure}
\else
  \newcommand\figurename{Figure}
\fi
\ifdefined\tablename
  \renewcommand*\tablename{Table}
\else
  \newcommand\tablename{Table}
\fi
}
\@ifpackageloaded{float}{}{\usepackage{float}}
\floatstyle{ruled}
\@ifundefined{c@chapter}{\newfloat{codelisting}{h}{lop}}{\newfloat{codelisting}{h}{lop}[chapter]}
\floatname{codelisting}{Listing}
\newcommand*\listoflistings{\listof{codelisting}{List of Listings}}
\makeatother
\makeatletter
\makeatother
\makeatletter
\@ifpackageloaded{caption}{}{\usepackage{caption}}
\@ifpackageloaded{subcaption}{}{\usepackage{subcaption}}
\makeatother
\usepackage{bookmark}
\IfFileExists{xurl.sty}{\usepackage{xurl}}{} % add URL line breaks if available
\urlstyle{same}
\hypersetup{
  pdftitle={Placing Archival Maps},
  colorlinks=true,
  linkcolor={blue},
  filecolor={Maroon},
  citecolor={Blue},
  urlcolor={Blue},
  pdfcreator={LaTeX via pandoc}}


\title{Placing Archival Maps}
\usepackage{etoolbox}
\makeatletter
\providecommand{\subtitle}[1]{% add subtitle to \maketitle
  \apptocmd{\@title}{\par {\large #1 \par}}{}{}
}
\makeatother
\subtitle{Assignment 3}
\author{}
\date{}
\begin{document}
\maketitle


\begin{longtable}[]{@{}ll@{}}
\toprule\noalign{}
Assigned & Due \\
\midrule\noalign{}
\endhead
\bottomrule\noalign{}
\endlastfoot
March 9, 2026 & April 3, 2026 \\
\end{longtable}

\section{Assignment overview}\label{assignment-overview}

This group assignment asks you to create an ``exhibit'' or product that
situates the map(s) your group is focusing on within their historical
and geographic context. Rather than accept these maps as honest and
straightforward representations of the world, the exhibit/product you
create should highlight how they were created by people in particular
positions to achieve a historically and geographically specific set of
goals. This is not to say you need to pass judgment on these
goals--being critical does not always mean taking an antagonistic
position. Rather, your group should help us understand how the map(s)
you are focusing on reflect choices about what or who should be mapped,
whose interests should be prioritized, and how to best represent the
world cartographically.

\subsection{Let's make that more
concrete!}\label{lets-make-that-more-concrete}

I'm suggesting two main types of products for your group to consider.

\begin{enumerate}
\def\labelenumi{\arabic{enumi}.}
\item
  \textbf{A virtual exhibit using the ArcGIS Storymap platform}. With
  this approach, you would include map(s) from your main object along
  with 6-8 other items from special collections. That would be
  structured in the following way:

  \begin{enumerate}
  \def\labelenumii{\arabic{enumii}.}
  \item
    A couple of paragraphs providing an \textbf{introduction} to your
    exhibit. This can include a description of the map(s) that are the
    main focus, the central theme you want to emphasize about those
    maps, and anything else that a viewer should keep in mind in looking
    through the artifacts you've collected.
  \item
    \textbf{A} \textbf{photo/image, citation, and caption} for your main
    map(s) and each artifact in your exhibit. The photo can be from your
    phone, or you can use a digitized version if available. The citation
    should be in APA
    format--\href{https://libguides.csudh.edu/archives-citation/apa}{here's
    one guide} to using that for archival materials. We will discuss
    captions in class, but they should be relatively short (no more than
    a paragraph) and provide notable details/features that viewers
    should pay attention to and an explanation of how this artifact fits
    within the broader exhibit.
    \href{https://drive.google.com/file/d/1QzaDQND_rpnHx2N8SHJnkhxC4VwiNM2T/view?usp=sharing}{Here's
    one chapter} on writing good captions.
  \item
    Lastly, you should have a final section of a few paragraphs that
    situates the exhibit you've created within \textbf{the context of
    our class}. What are the connections to our readings thus far in the
    course? How does this exhibit help sharpen our thinking about the
    role of maps in urban spaces and urban life? Think of it as an
    explanation of how your exhibit fits within the museum (our class).~
  \end{enumerate}
\item
  \textbf{A visual mosaic of map(s) from your main item along with other
  relevant archival materials.} With this approach, think of the
  conversation we had with Meghan Kelly and the work of Margaret Pearce
  on maps like
  ``\href{https://storymaps.arcgis.com/stories/2b5792be16d64f9e999ec93cd436bc17}{They
  Would Not Take Me There}''. Or you could think of work
  \href{https://hyperallergic.com/remagining-monuments-to-make-them-resonate-locally-and-personally/}{to
  reinterpret historical monuments such as confederate statues}here in
  the Southeast. The goal for those cartographers is using visuals and
  text to situate maps, to understand the positionalities of mappers and
  those being mapped. The goal is to show how the map itself is not an
  objective mirror of reality but a historically and culturally specific
  visual object. It's to call attention to aspects of these maps we
  might otherwise take for granted. More concretely, you can either
  digitally or physically alter your map while also adding archival
  materials or other text/visuals to help contextualize it. Here, too,
  you should think about a story you want to tell about your map--how it
  shows the world in ways that reflect its intended purpose, time, and
  place. Your group can consider what this might look like. I've got
  funds to print things out if you want.\\
  \strut \\
  If you choose option 2, I'd also ask your group to provide an
  \textbf{\emph{interpretative guide}} of no more than 2 single spaced
  pages that provides an introduction to your visual (what's its title?,
  what's the main focus?) as well as walking us through your design. It
  should also cite any items your group used (see information on how to
  do that in choice 1 above)
\end{enumerate}

\section{Assignment criteria}\label{assignment-criteria}

For either option, I will be assessing your group's project based on the
following criteria:

\begin{itemize}
\item
  Your project should demonstrate your ability to use archival research
  techniques to locate and interpret materials in UGA's Special
  Collections library.
\item
  Your project should draw upon techniques, approaches, and concepts
  discussed in our class to this point.
\item
  Your project should provide insights into the cultural, political, and
  historical context for your map(s)
\item
  Your project should draw clear and thoughtful connections between the
  archival materials you include.
\end{itemize}

I'll be rating all of these criteria on the same scale as other
assignments (Definitely/good enough/not yet/not addressed), and your
group will have a chance to revise if the assignment is rated
incomplete.~

\section{Submission details}\label{submission-details}

This document can be submitted on ELC under Assignments. This can be a
link to your group's storymap or a digital image file along with your
interpretive guide

\section{Group topics}\label{group-topics}

\begin{itemize}
\item
  \href{https://galileo-uga.primo.exlibrisgroup.com/discovery/fulldisplay?docid=alma9910300523902959&context=L&vid=01GALI_UGA:UGA&lang=en&search_scope=MyInstitution&adaptor=Local\%20Search\%20Engine&tab=LibraryCatalog&query=any,contains,moving\%20ahead\%20with\%20macon\%20&sortby=rank&facet=library,include,2959\%E2\%80\%93530589490002959}{Moving
  Ahead With Macon}

  \begin{itemize}
  \item
    \href{https://github.com/jshannon75/geog4636_mapsandthecity/blob/main/items/MovingAheadWithMacon_map1.pdf}{Scanned
    map 1}
  \item
    \href{https://github.com/jshannon75/geog4636_mapsandthecity/blob/main/items/MovingAheadWithMacon_map2.pdf}{Scanned
    map 2}
  \end{itemize}
\item
  \href{https://dlg.usg.edu/record/gsu_planatl_22866?canvas=22&x=555&y=755&w=2927}{Techwood
  neighborhood: Social base map survey of Atlanta, GA}
\item
  \href{https://sclfind.libs.uga.edu/catalog/ms4069_aspace_ref83_bqh}{Metropolitan
  North Georgia Water Planning District records: Nutrient Loading in the
  Lower Chattahoochee}
\item
  \href{https://sclfind.libs.uga.edu/catalog/RBRL180RMC_aspace_ref209_no2}{Rodney
  Mims Cook papers on Urban Renewal in University Center: Box 10}
\item
  \href{https://sclfind.libs.uga.edu/catalog/ms3995_aspace_ref325_i9w}{Georgia
  Wildlife Federation: Alcovy River Greenway project}
\item
  \href{https://dlg.usg.edu/record/guan_1633_014c-001?canvas=0&x=400&y=400&w=1986}{Athens
  R-51 Urban Renewal Project}
\item
  \href{https://sclfind.libs.uga.edu/catalog/ms2901_aspace_ref2200_y3m}{Union,
  South Carolina: James Hunt plans for low rent housing}
\end{itemize}




\end{document}
